\begin{theorem}\label{Thm:AlmostGraphAut}
Assume that a group $\calG$ acts transitively on $V$ from the left, no index two subgroup of $\calG$ acts transitively on $V$, and for any $g\in \calG$, the set 
$
\widetilde \calN (\widetilde \calN(g(\scrA)))$ 
is contained in 
$g(\widetilde \calN(\widetilde \calN(\scrA)))
$. Let $\mu> 0$ be such that for any subgroup $H$ of $\calG$ of index two and for any $u, v\in V$ lying in the same $H$-orbit with $u\in \widetilde \calN(v)$, the inequalities 
\begin{equation}
\label{Eqn:muBdd}
|\theta_i(Hv)\cap Hu|
\geq \frac{|V|}{2\mu},
\quad
|\theta_i(H^c v)\cap H^c u|
\geq \frac{|V|}{2\mu}
\end{equation}
hold for some integer $i_{u,v}$ depending on $u, v$ and lying between $1$ and $d$. 
\begin{enumerate}
\item 
If 
\begin{equation}
\label{Eqn:V108bdd20}
\mu = 1, \quad 
|V| 
 \geq 
108
\left( 2 + \frac {(\nu-1)(2\nu -1)} {30d}  \right)  
\left( \frac{ds}{\varepsilon^2} + \frac{dr}\varepsilon +\frac{ds}{\varepsilon}\right) + s ,
\end{equation}
then the vertex-Cheeger constant of the graph $(V, E)^2$ is greater than $\frac{\varepsilon^{2}}{30d}$. 
\item 
If 
\begin{equation}
\label{Eqn:V108bdd}
\mu = 2, \quad 
|V| 
 \geq 
108
\left( 2 + \frac {(\nu-1)(2\nu -1)}  {72d} \right)  
\left( \frac{ds}{\varepsilon^2} + \frac{dr}\varepsilon +\frac{ds}{\varepsilon}\right) + s ,
\end{equation}
then the vertex-Cheeger constant of the graph $(V, E)^2$ is greater than $\frac{\varepsilon^{2}}{72d}$. 
\item 
If
\begin{equation}
\label{Eqn:V108bdd20d3}
\mu = d, \quad 
|V| 
 \geq 
108
\left( 2 + \frac {(\nu-1)(2\nu -1)}  {30d^3} \right)  
\left( \frac{d^3s}{\varepsilon^2} + \frac{d^3r}\varepsilon+ \frac{d^{3}s}{\varepsilon}\right) + s,
\end{equation}
then the vertex-Cheeger constant of the graph $(V, E)^2$ is greater than $ \frac{\varepsilon^{2}}{30d^3}$. 
\end{enumerate}
\end{theorem}

\begin{proof}
On the contrary, assume that the vertex-Cheeger constant $\tth$ of the graph $(V, E)^2$ satisfies $\tth \leq  \frac{\varepsilon^{2}}{30d}$. 
Set 
\begin{align*}
r & =  1- \frac {d \tth}{ \varepsilon^2}  ( \varepsilon + \nu + 2).
\end{align*}
Since $\varepsilon \leq \nu -1$, it follows that 
$$
1 - r = \frac {d \tth}{ \varepsilon^2}  ( \varepsilon + \nu + 2)
\leq \frac {d \tth}{ \varepsilon^2}  (2\nu+1) 
\leq \frac{2\nu +1}{30}
\leq \frac 3{10}.$$
Using $\tth \leq  \frac{\varepsilon^{2}}{30d}$ and Equation~\eqref{Eqn:V108bdd20}, we obtain 
\begin{align*}
|\scrA| \left( 2 + \tth + \frac {\nu \tth} {\varepsilon} \right) 
& \geq 
|V| - s\\
& \geq 
108 \left( 2 + \frac {(\nu-1)(2\nu -1)} {30d}  \right)  
\left( \frac{ds}{\varepsilon^2} + \frac{dr}\varepsilon + \frac{ds}{\varepsilon}\right) \\
& \geq 
108 \left( 2 + \frac {\varepsilon (\varepsilon  + \nu)} {30d}  \right)  
\left( \frac{ds}{\varepsilon^2} + \frac{dr}\varepsilon + \frac{ds}{\varepsilon}\right) \\
& \geq 
108 \left( 2 + \tth + \frac {\nu\tth} {\varepsilon} \right)  
\left( \frac{ds}{\varepsilon^2} + \frac{dr}\varepsilon + \frac{ds}{\varepsilon}\right) ,
\end{align*}
which implies 
$ \frac{ ds}{\varepsilon^2 }
+ \frac{ dr}{\varepsilon } + \frac{ds}{\varepsilon}\leq \frac{|\scrA|}{108}$. 
Similarly, if $\tth \leq  \frac{\varepsilon^{2}}{72d}$, then Equation~\eqref{Eqn:V108bdd} implies 
$ \frac{ ds}{\varepsilon^2 }
+ \frac{ dr}{\varepsilon } + \frac{ds}{\varepsilon} \leq \frac{|\scrA|}{108}$. 
Moreover, if $\tth \leq  \frac{\varepsilon^{2}}{30d^3}$, then Equation~\eqref{Eqn:V108bdd20d3} implies 
$ \frac{ ds}{\varepsilon^2 }
+ \frac{ dr}{\varepsilon } + \frac{ds}{\varepsilon}\leq \frac{|\scrA|}{108}$. Consequently, $r\geq \frac 7{10}$ and 
$$
\frac {2d\tth} { \varepsilon^2} ( 2\nu + 1) 
+ \frac{2d s}{\varepsilon^2 |\scrA|}
+ \frac{3dr}{\varepsilon |\scrA|}
+ \frac{2ds}{\varepsilon |\scrA|}
\leq 
\frac 3{5} + \frac 3 {108} 
< 1.
$$
 
Consider the subset $H$ of $\calG$ defined by
$$H  :=
\left\{
g\in \calG \,:\,\, |\scrA \cap  (g(\scrA))| \geq r|\scrA| - \frac{ds}{\varepsilon^2} 
- \frac{dr}{\varepsilon}
- \frac{ds}{\varepsilon}
\right\}.$$
Note that $H$ contains the identity element of $\calG$. Let $g_1, g_2$ be elements of $H$. By the triangle inequality, 
\begin{align*}
|\scrA \setminus (g_1(g_2(\scrA)))| 
& \leq |\scrA \setminus (g_1(\scrA)) | + |(g_1(\scrA)) \setminus (g_1(g_2(\scrA)))| \\
& =  |\scrA \setminus  (g_1(\scrA)) | + |\scrA  \setminus (g_2(\scrA)) |\\
& = |\scrA | -  |\scrA \cap (g_1(\scrA))  | + |\scrA | - |\scrA  \cap (g_2(\scrA)) |\\
& \leq 2|\scrA| - 2r |\scrA| + \frac{2ds}{\varepsilon^2}  + \frac{2dr}{\varepsilon} 
+\frac{2ds}{\varepsilon} ,
\end{align*}
which yields that  
\begin{align*}
|\scrA \cap (g_1(g_2(\scrA)))| 
& = |\scrA | - |\scrA \setminus (g_1(g_2(\scrA)))| \\
& \geq |\scrA | - 2|\scrA| + 2r |\scrA| - \left(\frac{2ds}{\varepsilon^2} +\frac{2dr}{\varepsilon} + \frac{2ds}{\varepsilon} \right)\\
& = (2r - 1) |\scrA| - \left(\frac{2ds}{\varepsilon^2} + \frac{2dr}{\varepsilon} +\frac{2ds}{\varepsilon} \right).
\end{align*}
If $|\scrA \cap (g_1(g_2(\scrA)))| \leq (1-r) |\scrA| + \frac{ds}{\varepsilon^2} + \frac{2dr}\varepsilon +\frac{ds}{\varepsilon}$, then we obtain 
\begin{align*}
(1-r) |\scrA| + \frac{ds}{\varepsilon^2} + \frac{2dr}\varepsilon + \frac{ds}{\varepsilon}
&\geq |\scrA \cap (g_1(g_2(\scrA)))| \\
&\geq  (2r - 1) |\scrA| - \left(\frac{2ds}{\varepsilon^2} + \frac{2dr}{\varepsilon} + \frac{2ds}{\varepsilon} \right),
\end{align*}
which yields
$(3r-2)
|\scrA| 
\leq \frac{3ds}{\varepsilon^2} + \frac{4dr}{\varepsilon} + \frac{3ds}{\varepsilon} 
\leq \frac {|\scrA|}{ 11},
$
which contradicts $ r \geq \frac 7{10}$. 
So, the inequality 
$|\scrA \cap (g_1(g_2(\scrA)))| \leq (1-r) |\scrA|
+ \frac{ds}{\varepsilon^2} + \frac{2dr}\varepsilon + \frac{ds}{\varepsilon}$
does not hold. By Proposition~\ref{Prop:AExists}(2), $H$ contains the product $g_1g_2$. Consequently, $H$ is a subgroup of $\calG$. 

Let $t$ denote the cardinality of the stabiliser of an element of $V$ under the action of $\calG$. Note that for any $g\in \calG$, the map
$$\scrA\cap g^\mo \scrA \to \scrA\times \scrA, \quad 
x\mapsto (x, gx)$$ 
induces a map 
$$\varphi: \coprod _{g\in \calG} \scrA\cap g^\mo \scrA \to \scrA\times \scrA.$$
Each fibre of this map contains exactly $t$ elements. This implies 
$$
|\scrA|^2
= |im(\varphi)|
= \sum_{x\in im(\varphi) } \frac{|\varphi^\mo (x)|} {|\varphi^\mo (x)|}
= \frac 1t \sum_{x\in im(\varphi)} |\varphi^\mo (x)|
= \frac 1t \sum_{g\in \calG} |\scrA\cap g^\mo \scrA|.$$

If $\calG = H$, then 
$$|\scrA| \cdot \frac{|V|}{2} 
\geq |\scrA|^2 = \frac 1t \sum_{g\in \calG} |\scrA\cap g^\mo \scrA|
\geq \frac 1t |\calG| \cdot 
\left(r|\scrA|  - \frac{ds}{\varepsilon^2} 
- \frac{dr}{\varepsilon}
- \frac{ds}{\varepsilon}
\right),$$
which implies $r \leq \frac 12 + \frac 1{108}$. This contradicts the bound $r \geq \frac 7 {10}$. So, $H$ is a proper subgroup of $\calG$. 

The estimate 
\begin{align*}
t |\scrA|^2 
&= \sum_{g\in \calG} |\scrA\cap g^\mo \scrA| \\
&\leq |H| |\scrA| + \left(
\frac{d\tth}{\varepsilon^2}(\varepsilon + \nu+ 2)|\scrA|
+ \frac{ds}{\varepsilon^2} + \frac{2dr}\varepsilon + \frac{ds}{\varepsilon}
\right)
|\calG\setminus H|
\end{align*}
yields
$$t |\scrA|  
\leq |H| + 
\left(
\frac{d\tth}{\varepsilon^2}(\varepsilon + \nu+ 2)
+ 
\alpha
\right)
(|\calG| - |H|)$$
with $\alpha  = \frac{ds}{\varepsilon^2|\scrA|} + \frac{2dr}{\varepsilon |\scrA|} + \frac{ds}{\varepsilon |\scrA|}$. 
Proposition~\ref{Prop:AExists}(1) implies
\begin{align*}
\left(\frac{1}{2+ \tth + \frac{\nu\tth}{\varepsilon} +\frac{s}{|\scrA|}}\right) |\calG| 
& = \left(\frac{t}{2+ \tth + \frac{\nu\tth}{\varepsilon} + \frac{s}{|\scrA|}}\right) |V| \\
&\leq 
 \left(
\frac{d\tth}{\varepsilon^2}(\varepsilon + \nu+ 2)
+ 
\alpha
\right)|\calG| 
+  \left(1 - \frac{d\tth}{\varepsilon^2}(\varepsilon + \nu+ 2)
- \alpha
\right) |H|.
\end{align*}
Now we show that $H$ is a subgroup of $\calG$ of index two. To establish this fact, we use 
\[\frac 13 \left(1 - \frac{d\tth}{\varepsilon^2}(\varepsilon + \nu+ 2) - \alpha \right)
< 
\left(\frac{1}{2+ \tth + \frac{\nu\tth}{\varepsilon}+ \frac{s}{|\scrA|}}\right) - \frac{d\tth}{\varepsilon^2}(\varepsilon + \nu+ 2) - \alpha,\]
equivalently, 
$$\left(2 + \tth + \frac{\nu\tth}{\varepsilon} +\frac{s}{|\scrA|}\right) \left( 1 + \frac{2d\tth}{\varepsilon^2}(\varepsilon + \nu+ 2) + 2\alpha \right) < 3.$$
Since $\tth\leq  \frac{\varepsilon^{2}}{2x d}$ holds with $x= 15$, it follows that  
\begin{align*}
& \left(2 + \tth + \frac{\nu\tth}{\varepsilon} + \frac{s}{|\scrA|}\right) \left( 1 + \frac{2d\tth}{\varepsilon^2}(\varepsilon + \nu+ 2) +
2\alpha \right)\\
& 
\leq 
\left(2 + \frac{\tth}{\varepsilon} (\nu + \varepsilon) + \frac{s}{|\scrA|}\right) \left( 1 + \frac{2d\tth}{\varepsilon^2}(\varepsilon + \nu+ 2) + \frac{2ds}{\varepsilon^2|\scrA|} + \frac{4dr}{\varepsilon |\scrA| } + \frac{2ds}{\varepsilon |\scrA|} \right)\\
& 
\leq 
\left(2 + \frac{\tth}{\varepsilon} (2\nu -1) + \frac{s}{|\scrA|} \right) \left( 1 + \frac{2d\tth}{\varepsilon^2}(2\nu+ 1) + \frac{2ds}{\varepsilon^2|\scrA|} + \frac{4dr}{\varepsilon |\scrA|} + \frac{2ds}{\varepsilon |\scrA|}  \right)\\
& 
< 
\left(2 + \frac{(\nu-1)(2\nu -1)}{2xd} +\frac{1}{108} \right) \left( 1 + \frac{2\nu+ 1}{x}+ \frac 1{27}\right)\\
& 
\leq 
\begin{cases}
\left(2 + \frac 3{4x}+ \frac{1}{108}\right) \left( 1 + \frac{5}{x}+ \frac 1{27}\right) & \text{ if } d = 2, \\
\left(2 + \frac 1{x} + \frac{1}{108}\right) \left( 1 + \frac{6}{x}+ \frac 1{27}\right) & \text{ if } d \geq 3,
\end{cases}\\
& 
< 3.
\end{align*}
This proves the claim that $H$ is a subgroup of $\calG$ of index two. Since no index two subgroup of $\calG$ acts transitively on $V$, the action of $H$ on $V$ has at least two distinct orbits. Since the action of $\calG$ on $V$ is transitive, the action of $H$ on $V$ has exactly two orbits. Denote the orbits of the action of $H$ on $V$ by $\calO_1, \calO_2$.  

Now we use Fre\u{\i}man's technique~\cite{FreimanGroupsInverseProb} to conclude the proof. For any $g\in G \setminus H$ and for any subsets $B, C$ of $V$ contained in $\calO_1$, $\calO_2$ (in some order), the map 
$B\cap g^\mo C \to B\times C$, sending 
$x\mapsto (x, gx)$, induces a map 
$\varphi: \coprod _{g\in H^c} B\cap g^\mo C \to B\times C$. 
If $B \times C$ is nonempty, then each fibre of this map contains exactly $t$ elements and hence 
$$
|B\times C|
= |im(\varphi)|
= \sum_{x\in im(\varphi) } \frac{|\varphi^\mo (x)|} {|\varphi^\mo (x)|}
= \frac 1 t \sum_{x\in im(\varphi)} |\varphi^\mo (x)|
= \frac 1 t \sum_{g\in H^c} |B\cap g^\mo C|.$$
Note that 
$|B\times C|
= \frac 1 t \sum_{g\in H^c} |B\cap g^\mo C|$
holds even if $B \times C = \emptyset$. 
For any $g\in H^c$, we have 
\begin{align*}
\scrA\cap g^\mo \scrA
& = 
((\scrA\cap \calO_1) \cup (\scrA\cap \calO_2))
\cap 
(g^\mo ((\scrA\cap \calO_1) \cup (\scrA\cap \calO_2)))\\
& = 
((\scrA\cap \calO_1) \cap 
(g^\mo (\scrA\cap \calO_2))
\cup 
((\scrA\cap \calO_2) \cap 
(g^\mo (\scrA\cap \calO_1)),
\end{align*}
which implies 
\begin{align*}
& |\scrA \cap \calO_1| |\scrA\cap \calO_2|\\
& = \frac 12 |\scrA \cap \calO_1| |\scrA\cap \calO_2| + 
\frac 12 |\scrA \cap \calO_2| |\scrA\cap \calO_1|\\
& =
\frac 1 {2t} \sum_{g\in H^c} |(\scrA \cap \calO_1) \cap g^\mo (\scrA \cap \calO_2)| 
+ 
\frac 1{2 t} \sum_{g\in H^c} |(\scrA \cap \calO_2) \cap g^\mo (\scrA \cap \calO_1)| \\
& = 
\frac 1 {2t} \sum_{g\in H^c} |\scrA \cap g^\mo \scrA|\\
& 
\leq \frac 1 {2t} 
\left(
\frac{d \tth}{\varepsilon^2} (\varepsilon + \nu + 2)|\scrA| 
+ \frac{ds}{\varepsilon^2} + \frac{2dr}\varepsilon + \frac{ds}{\varepsilon}
\right)|H|\\
& 
= \frac 1 {2t} 
\left(
\frac{d \tth}{\varepsilon^2} (\varepsilon + \nu + 2)|\scrA| 
+\frac{ds}{\varepsilon^2} + \frac{2dr}\varepsilon + \frac{ds}{\varepsilon}
\right)
\frac{t|V|}2\\
& 
\leq 
\frac{d \tth}{\varepsilon^2} (\varepsilon + \nu + 2)
\frac{|V|^2}8
+ 
\left(\frac{ds}{\varepsilon^2} + \frac{2dr}\varepsilon + \frac{ds}{\varepsilon}
\right)
\frac{|V|}4.
\end{align*}
So, the orbit $\calO$ of some element of $V$ under the action of $H$ satisfies 
$$
|\scrA\cap \calO^c| \leq 
\sqrt{\frac{d\tth}{8\varepsilon^2} (\varepsilon + \nu + 2) 
+
\left(\frac{ds}{\varepsilon^2} + \frac{2dr}\varepsilon + \frac{ds}{\varepsilon}
\right)
\frac{1}{4|V|}
}|V|.$$
If the inequalities $\tth<  \frac{\varepsilon^{2}}{2x d^{1 + 2\delta}}$ and $ \frac{ d^{1 + 2\delta} s}{\varepsilon^2 }
+ \frac{ d^{1 + 2\delta} r}{\varepsilon } + \frac{d^{1 + 2\delta} s}{\varepsilon}\leq \frac{|\scrA|}{108}$ hold with $\delta \in \{0, 1\}$, then for any $1\leq i \leq d$, 
\begin{align*}
&|\theta_i(\calO)\cap \calO |\\
& = |\theta_i(\calO\cap \scrA) \cap \calO| + |\theta_i(\calO\setminus \scrA) \cap \calO|\\
& \leq |\theta_i(\scrA) \cap \calO| + |\calO\setminus \scrA|\\
& = |\theta_i(\scrA) \cap (\calO\cap \scrA) | + |\theta_i(\calO\cap \scrA) \cap (\calO\setminus \scrA) | + |\calO\setminus \scrA|\\ 
& \leq |\theta_i(\scrA) \cap \scrA| + 2|\calO\setminus \scrA|\\ 
& = |\theta_i(\scrA) \cap \scrA| + 2|\calO| - 2|\calO\cap \scrA|\\ 
& = |\theta_i(\scrA) \cap \scrA| + |V| - 2|\scrA| + 2|\scrA\cap \calO^c|\\ 
& \leq s +  |\calN(\scrA) \cap \scrA| + s + \left(2 + \tth + \frac{\nu\tth}{\varepsilon}\right)|\scrA|- 2|\scrA| + 2|\scrA\cap \calO^c|\\ 
& \leq 2s + \frac \tth\varepsilon |\scrA| + \left( \tth + \frac{\nu\tth}{\varepsilon}\right)|\scrA| + 2\sqrt{\frac{d\tth}{8\varepsilon^2} (\varepsilon + \nu + 2) 
	+
	\left(\frac{ds}{\varepsilon^2} + \frac{2dr}\varepsilon + \frac{ds}{\varepsilon}
	\right)
	\frac{1}{4|V|}
}|V| \\
& \leq 2s + \left(\frac \tth\varepsilon (\nu+1 + \varepsilon)  + \sqrt{\frac{2d\tth}{\varepsilon^2} (\varepsilon + \nu + 2) 
+\left(\frac{ds}{\varepsilon^2} + \frac{2dr}\varepsilon + \frac{ds}{\varepsilon}
\right)
\frac{4}{|V|}
}\right) \frac{|V|}{2} \\
& \leq 2s + \left(\frac {2\nu\tth}\varepsilon + \sqrt{\frac{2d\tth}{\varepsilon^2} (2\nu + 1) 
+\left(\frac{ds}{\varepsilon^2} + \frac{2dr}\varepsilon + \frac{ds}{\varepsilon}
\right)
\frac{4}{|V|}
}\right) \frac{|V|}{2} \\
& \leq \frac{\varepsilon^2}{d^{1 + 2\delta}} \frac{|\scrA|}{54} + \left(\frac {\nu\varepsilon}{xd^{1 + 2\delta}} + \frac 1 {d^{\delta}}\sqrt{\frac{ 2\nu + 1 }{x}
+
\frac{4}{108}
}\right) \frac{|V|}{2} \\
& \leq \frac{(\nu -1)^2}{d^{1 + 2\delta}} \frac{|V|}{108} + \left(\frac {\nu(\nu -1)}{xd^{1 + 2\delta}} + \frac 1 {d^{\delta}}\sqrt{\frac{ 2\nu + 1 }{x}
+
\frac{4}{108}
}\right) \frac{|V|}{2} \\
& \leq \left(\frac{(\nu -1)^2}{54 d^{1 + 2\delta}} + \frac {\nu(\nu -1)}{xd^{1 + 2\delta}} + \frac 1 {d^{\delta}}\sqrt{\frac{ 2\nu + 1 }{x}
+
\frac{4}{108}
}\right) \frac{|V|}{2} \\
& \leq \left(\frac{(\nu -1)^2}{54 d} + \frac {\nu(\nu -1)}{xd} + \sqrt{\frac{ 2\nu + 1 }{x}
+
\frac{4}{108}
}\right) \frac{|V|}{2d^{\delta}} \\
& 
\leq 
\begin{cases}
\left(\frac 1{108}+ \frac {1}{x} + \sqrt{\frac{ 5 }{x}
+
\frac{4}{108}
}\right) \frac{|V|}{2 d^{\delta}} 
& \text{ if } d =2, \\
\left(\frac 1{72} + \frac {5}{4x} + \sqrt{\frac{6}{x}
+
\frac{4}{108}
}\right) \frac{|V|}{2 d^{\delta}} 
& \text{ if } d \geq 3, 
\end{cases}
\\
& 
<
\begin{cases}
 \frac{|V|}{2d^{\delta}} 
& \text{ if } x \geq 9, \\
\frac{|V|}{4d^{\delta}} 
& \text{ if } x \geq 36.
\end{cases}
\end{align*}
So, under the hypothesis of each of statements (1), (2), (3), it follows that 
$$
|\theta_i(\calO)\cap \calO|
< \frac{|V|}{2\mu}
$$
for any $1\leq i \leq d$. Suppose two vertices $u, v$ of $\calO^\dag$ are adjacent in $(V, \widetilde E)$ where $\calO^\dag$ is one of $\calO, \calO^c = V \setminus \calO$. Let $H^\dag$ denote $H$ (resp. $H^c = \calG \setminus H$) if $\calO^\dag  = \calO$ (resp.~$\calO^\dag = \calO^c$). Note that $\calO = H^\dag u = H^\dag v$. Using Equation~\eqref{Eqn:muBdd}, we obtain that 
$
|\theta_i(\calO)\cap \calO|
 = |\theta_i(H^\dag v) \cap H^\dag u | 
 \geq \frac{|V|}{2\mu}
$
for some $1\leq i \leq d$. It follows that $\calO$ and $\calO^c$ are independent subsets of $(V, \widetilde E)$, implying that $(V, \widetilde E)$ is bipartite, which in turn shows that $(V, E)$ is bipartite, contradicting the hypothesis. So, the vertex-Cheeger constant of the graph $(V, E)^2$ satisfies the claimed bounds. 
\end{proof}
